\subsection{软件系统与模块设计}
软件分为离线训练程序和板上控制程序。训练程序将十组参数导出为整数参数包，随配置文件进入板上ROM。板上的Sapphire RV32IM运行C固件，在32\,KiB片内RAM中维护当前风格、视图请求和手势模板。C程序经APB读写设备寄存器，像素则经AXI直接在DDR与专用RTL模块之间传输。这一分工使处理器可以在卷积计算期间继续响应按键和分析手势快照。

\begin{table}[H]\centering\small
\caption{板上视频软件与专用逻辑的模块接口}\label{tab:swroles}
\begin{tabularx}{\linewidth}{@{}p{.23\linewidth}p{.32\linewidth}X@{}}\toprule
模块、执行位置 & 主要输入 & 输出及下游使用方式\\\midrule
设备与风格控制（C） & 目标风格、采集及引擎状态、ROM参数 & I\textsuperscript{2}C配置，APB参数地址/数据，作业启动。PLIC IRQ16返回作业事件。\\
采集与源帧管理（RTL） & 四lane PPI字节及状态，RAW行与帧边界 & RGB565源帧。给作业控制器提供最新完整帧号和可锁定的槽。\\
分条DMA与卷积核（RTL） & 源、目的槽地址，已加载权重，本次视图 & 结果帧及完成/错误信号。完成条件包含最后一笔DDR写响应。\\
显示读取与叠加（RTL） & 待提交槽号、帧号、视图和风格，C/I等状态 & 在消隐期确认切换，输出同步、数据有效和RGB，驱动TMDS。\\
手势交互（RTL＋C） & 有效原图区域、模式、按键，带序号的64×64快照 & C程序输出候选与确认进度，确认后向控制器提交目标风格。\\\bottomrule
\end{tabularx}\end{table}

图\ref{fig:frame_messages}以一次风格作业说明控制关系。软件可以在收到计算完成中断后调度下一幅图，而一幅图究竟何时出现在屏上，由显示端的消隐确认决定。因此屏上的风格编号取自前台帧的元数据。
\begin{figure}[H]\centering
\begin{tikzpicture}[x=1cm,y=.72cm,font=\small,
lane/.style={draw=navy!65,fill=pale,rounded corners=2pt,minimum width=2.55cm,minimum height=.58cm},
msg/.style={-{Stealth[length=2mm]},draw=teal,thick},
lbl/.style={midway,above,inner sep=1pt,font=\footnotesize}]
\foreach \x/\t in {0/RISC-V固件,3.55/帧与作业控制,7.1/CNN及DMA,10.65/显示读取} {
 \node[lane] at (\x,0){\t};\draw[navy!35,dashed] (\x,-.45)--(\x,-6.3);}
\draw[msg] (0,-.95)--node[lbl]{参数就绪，启动作业}(3.55,-.95);
\draw[msg] (3.55,-1.95)--node[lbl]{锁住源帧，指定输出槽}(7.1,-1.95);
\draw[msg] (7.1,-2.95)--node[lbl]{计算结束，最后B已返回}(3.55,-2.95);
\draw[msg] (3.55,-3.9)--node[lbl]{完成事件经PLIC通知}(0,-3.9);
\draw[msg] (3.55,-4.8)--node[lbl]{提交完整结果及帧号}(10.65,-4.8);
\draw[msg] (10.65,-5.8)--node[lbl]{安全消隐确认，释放旧前台}(3.55,-5.8);
\end{tikzpicture}
\caption{一次风格作业的关键事件。从上往下阅读。中断和显示提交使用各自的握手，图中行距不表示时间间隔。}\label{fig:frame_messages}
\end{figure}

\subsubsection{采集与源帧管理}
固件释放模组复位，通过I\textsuperscript{2}C读取寄存器\sig{0x3107/0x3108}，确认器件ID为\sig{0xCD6B}，设置并读回1920×1080输出尺寸，再根据实际帧长限制曝光。摄像头开流后，D-PHY输出的四路字节、有效位和同步状态被打成一条44位记录。深度为512的异步FIFO将整条记录从接收字节时钟域送入100\,MHz的用户接口域。即使摄像头在低功耗间隔暂停字节时钟，同一批字节与其状态仍能对应。

CSI解析模块核验包头ECC、行负载长度及CRC。1920个RAW10样本每行占2400字节，每五字节还原四个十位值。Bayer模块取样本高八位，以两条行缓存形成$3\times3$邻域。例如在蓝色像素处，中心值保留为B，四个正交邻点的平均值为G，四个对角邻点的平均值为R。缩放逻辑沿横纵两轴按$0,1,3,4,\ldots$选取插值后的像素，得到1280×720图像。每个有效输出拍含两个RGB565像素，较左像素占低16位。

\sig{camera_pack_720}将四拍双像素组成一个128位DDR词，并携带源帧号和17位字索引。行跨度为2560字节，故像素$(x,y)$的地址为$A_{\rm src}+2560y+2x$。一帧共115200个词，占1843200字节。帧尾记录、连续词数和最后的AXI写响应均有效时，源池才发布这幅图。格式错误或写响应错误会使整帧作废。

源池的三个槽从4、6、8\,MiB起始。\sig{source_frame_pool}在作业领取最新好帧时锁住其槽，采集端选择另一个空槽继续写。当新帧到来而上一幅完整帧仍未被领取，池中以新帧替换它并累计跳过数。该计数与摄像头传输坏帧分开，可用于区分采集质量和推理是否跟得上输入。

\subsubsection{风格调度与分条计算}
假设当前正显示梵高，操作者按键选择水墨。按键模块只更新请求槽号。固件在引擎空闲时，从APB基址\sig{0xF8100000}加上$\mathtt{0x4000}+736s$的位置读取槽$s$的参数，按层写入权重、偏置和移位。跨域邮箱保持地址和数据，直到166.667\,MHz的CNN域完成写入并应答。处理器核对216项配置后提交新的已加载槽号。当完整源帧可用时再启动作业。若中途又收到操作，当前作业使用启动时锁存的风格，后续作业处理最新请求。

整幅六通道特征图占$1280\times720\times6$字节，远大于处理一个局部窗口所需的存储。引擎沿宽度划分四条，每条交付中央320列，并在两侧各保留16列邻域供卷积读取。图\ref{fig:tilegeometry}以第二条为例。它输出全图第320--639列，输入窗口横跨304--655列，同一段边缘可以由相邻两条共享。宽度352列的行历史在四条之间复用。

\begin{figure}[H]\centering
\begin{tikzpicture}[x=1cm,y=1cm,font=\small]
\foreach \i/\a/\b in {0/0/319,1/320/639,2/640/959,3/960/1279}{
 \pgfmathsetmacro{\lx}{3*\i}
 \ifnum\i=1 \fill[teal!22] (\lx,0) rectangle ++(3,.57);\else\fill[pale] (\lx,0) rectangle ++(3,.57);\fi
 \draw[navy!65] (\lx,0) rectangle ++(3,.57);
 \node at (\lx+1.5,.285){\a--\b};}
\node[anchor=west,font=\footnotesize] at (0,.89){全图横坐标（列），每条均处理720行};
\draw[navy!55,dashed] (3,0)--(2.2,-.7);\draw[navy!55,dashed] (6,0)--(9.8,-.7);
\filldraw[fill=orange!16,draw=navy!70] (0,-1.35) rectangle (2.2,-.7);
\filldraw[fill=teal!22,draw=navy!70] (2.2,-1.35) rectangle (9.8,-.7);
\filldraw[fill=orange!16,draw=navy!70] (9.8,-1.35) rectangle (12,-.7);
\node at (1.1,-1.02){304--319};\node at (6,-1.02){320--639};\node at (10.9,-1.02){640--655};
\node[font=\footnotesize] at (1.1,-1.63){邻域16列};\node[font=\footnotesize] at (6,-1.63){写回DDR的320列};\node[font=\footnotesize] at (10.9,-1.63){邻域16列};
\end{tikzpicture}
\caption{四条带覆盖全幅图像。下栏展开第二条，邻域宽度为便于阅读而放大。四条合计写回921600个独立像素。}\label{fig:tilegeometry}
\end{figure}

感受野半径为15像素，16列重叠足以让中央区域取得完整邻域。各卷积层以全图边缘为镜像边界，分条接缝保留真实邻点，因而物体跨过接缝时仍有连续的笔触。DW用行缓存和三组列存储提供空间取样。PW用两条输出通道并行运算，由四组上下文交叠处理不同像素，结果按输入顺序送往下一层。原RGB经FIFO与头部增量对齐后相加，输入停顿或输出背压时两条支路保持对应。

分条引擎\sig{video_tile_engine_async}每次通过DMA突发读16个128位词，将256字节缓存到片上。卷积结果每八个像素拼成一个128位写词，两组写缓冲交替拼包和等待总线应答。必须等四条的有效列写完并收到最后一个写响应，才将结果加入输出队列。三个输出槽位于0、2、10\,MiB，可分别保存正在计算的帧、待提交的帧与正在扫描的前台帧。完成事件经PLIC通知固件，处理器可在显示端等待消隐期间调度下一幅图。

\subsubsection{显示合成与状态叠加}
原图视图把源帧复制到输出槽，风格视图写入卷积结果。左右对比在写回前完成合成。对每一个输出坐标，逻辑以全局$x<640$选择对齐后的原RGB，其余列选择风格结果，合成后的整帧共用一个源帧号。这样，即使镜头前的手在运动，中线两侧仍对应同一次采集。每个输出槽同时保存视图与风格编号。

显示读取模块\sig{frame_reader}通过AXI把前台数据送入像素FIFO。像素端以74.25\,MHz产生每行1650拍、每帧750行的扫描时序，有效区为1280×720。结果队列提交新槽时，读取模块排空旧地址的在途读请求，并等待FIFO复位在读写两个时钟域完成。它在有效区开始前保留一行余量，确认仍处于消隐窗口才应答换帧。错过窗口就等待下一次消隐，避免同一次屏幕扫描混入两幅画面。

信息层按像素坐标叠加文字和手势框。C统计一秒内采集模块发布的好帧，I统计一秒内完成的神经网络作业，H60对应固定的HDMI扫描设置。显示器重复扫描前台时，I保持实际作业计数。S随新前台一同更新，因此可以用它判断风格何时真正出现。KEY3只切换文字和青框的可见性。USB串口以CAM7记录输出帧数、作业周期和错误计数，供连续测量使用。

\subsubsection{按键与手势交互}
KEY0在原图、风格图、左原右风格三种视图间循环，KEY1和KEY2在十个风格槽间前后选择。四路低有效按键经过两级同步和约10\,ms消抖，每次按下仅产生一次操作。手势与按键共用风格请求入口，区别在于手势先经过图像分类和持续时间判定。

当请求视图与前台一致，且处于原图或对比模式时，系统允许手势采样。采样器从原图左侧640×640的青框区域取样。采样器在横纵方向每10个像素取一点，得到64×64 RGB565快照。只有4096个采样点及好帧帧尾均已到达，处理器才可读取。快照附带序号和界面状态标识，同一张快照只用于一次判定。此旁路占8\,KiB，可以在主视频保持720p时独立使用。

首次使用个人手势时，按住KEY1与KEY2约一秒进入学习，在空框内采三张背景，再为每类手势各录四张。处理器先补偿整体亮度变化，再对背景差分结果保留最大连通区域，用亮度的第5和第95百分位规范化手形。每张样本提取192维特征，包括64维梯度方向、64维局部灰度及64维轮廓密度。软件比较正向和镜像的模板距离，并以最优距离及其与次优类别的间隔拒绝歧义。模板存于RAM，重新配置后须重新录入。

\begin{table}[H]\centering\small
\caption{数字手势与风格槽的对应关系。数字0与未检测到手势分别编码。}
\begin{tabularx}{\linewidth}{@{}clXclX@{}}\toprule
槽 & 风格 & 手形 & 槽 & 风格 & 手形\\\midrule
0 & 梵高厚涂 & 五指弯曲围成O形 & 5 & 阿什维尔抽象 & 五指张开\\
1 & 浮世绘 & 食指 & 6 & 形体对比 & 拇指、小指\\
2 & 彩色剪纸 & 食指、中指 & 7 & 蒙德里安几何 & 拇指、食指、小指\\
3 & 复古漫画 & 中指、无名指、小指 & 8 & 彩色马赛克 & 拇指、食指\\
4 & 竹林水墨 & 除拇指外四指 & 9 & 雨夜厚涂 & 食指弯钩\\\bottomrule
\end{tabularx}\end{table}

运行时每约333\,ms抓取一幅新快照，三张中的多数结果形成一次决定。第一个有效决定开启计时，约在第0、1、2、3、4秒连续得到同一数字，且实际经过四个完整秒后才提交。候选变化、界面切换或快照过期都会清除累计进度。原图确认后自动进入风格图，对比视图确认后仍保留左右两侧。图\ref{fig:gesture-time}用数字2示意一次成功识别。
\begin{figure}[H]\centering
\begin{tikzpicture}[x=2.1cm,y=1cm,font=\small]
\draw[-{Stealth},navy] (-.35,0)--(4.95,0);
\foreach \i in {0,1,2,3,4}{
 \draw[navy] (\i,-.07)--(\i,.07);\fill[teal] (\i,0) circle(2.1pt);
 \node[above] at (\i,.12){G2};\node[below] at (\i,-.12){P\i};}
\node[anchor=west,font=\footnotesize] at (4.18,.42){提交槽2};
\node[font=\footnotesize] at (2,-.85){相邻决定约间隔1秒，每个决定使用三张新快照};
\end{tikzpicture}
\caption{手势的四秒确认。首个决定记为P0，第五个同类决定且经过四秒可提交。前台实际更新后显示S2。}\label{fig:gesture-time}
\end{figure}
